% ============================================================
% 片段：s01  对应原文 section1.tex 第 1--42 行；PDF 第 4--6 页
% 内容：§1 Introduction（引言）
% 译者：lead   审校状态：初稿
% ============================================================

% ==== 原文 PDF 第 4 页 ====
\section{引言}

\begin{keytakeaways}
  \begin{itemize}[leftmargin=1.4em,itemsep=2pt,topsep=2pt]
    \item \textbf{物理图像}：对称性的现代语言不再只谈“群作用在局域算符上”，而是谈拓扑算符及其作用于延展对象的方式。由此得到两个推广方向：一是把作用对象从点推广到 $p$ 维延展对象（高形式对称性），二是把群结构本身松弛为不可逆的、范畴化的结构（非可逆对称性）。
    \item \textbf{数学工具}：拓扑算符与拓扑缺陷、平坦联络、分类空间、群上同调与配边（cobordism）分类、以及高阶范畴论。
    \item \textbf{关键结论}：普通 $0$-形式对称性只是最简单的一档；高形式对称性在规范理论中按 $U(1)$ 的电/磁中心荷与离散挠部分分类，而 't Hooft 反常与 SPT 相由同一份拓扑数据（配边群）统一刻画。本讲义按“拓扑生成元 $\to$ 高形式对称性 $\to$ 反常与 SPT $\to$ 应用 $\to$ 高群对称性”的顺序展开。
  \end{itemize}
\end{keytakeaways}

对称性是物理学的一条基本指导原则。在对称变换下，物理理论保持不变，而物理要素（场、算符）则被组织成对称性的不同表示。对称性有两种不同的含义：
\begin{enumerate}
  \item 整体对称性：对称参数 $g$ 与时空无关。整体对称性通常被视为物理系统的“真正的对称性”。
  \item 局域对称性（规范对称性）：对称参数 $g$ 依赖于时空。这种局域对称性通常也称为规范对称性，它被解释为理论的一种冗余，而不是真正的物理对称性。\footnote{此外还存在介于两者之间的对称性，例如自由实标量场 $\int \ptl_\mu\phi\ptl^\mu\phi$ 情形下的“Galileon 对称性”\cite{Nicolis:2008in,Hinterbichler:2015pqa} $\phi\rightarrow\phi+a$。其参数 $a$ 可以只要求满足 $\ptl_\mu\ptl^\mu a=0$。这是一个没有规范场的“半局域”对称性的例子。}
\end{enumerate}

传统上，“普通对称性”指的是对称群\footnote{群是一个集合，其上带有封闭、满足结合律的二元运算（记作 $\cdot$），在该运算下有单位元，且所有元素在该运算下可逆。}作用在局域算符\footnote{即只依赖于时空流形上一点的算符。}上的情形。在现代语言中，这被称为可逆的 $0$-形式对称性。

普通对称性（“普通”一词的含义后文会解释）的一个简单例子是

\begin{example}
\[ \mathcal{L} = |\partial_\mu \phi|^2 \]
其中 $\phi$ 是复标量场，其对称变换为 $\phi\mapsto e^{i\alpha}\phi$，$e^{i\alpha}\in U(1)$。
\end{example}

% ==== 原文 PDF 第 5 页 ====
过去十年间，人们从凝聚态物理、高能物理与纯数学等不同视角出发，为推广对称性的概念付出了大量努力\cite{Gaiotto:2014kfa}，参见以下综述与讲义\cite{Kong_2020,mcgreevy2022generalized,Cordova:2022ruw,Gomes:2023ahz,Schafer-Nameki:2023jdn,Brennan:2023mmt,Bhardwaj:2023kri}。从普通对称性的定义出发，可以观察到若干推广的线索：
\begin{enumerate}
  \item 从 $0$-形式对称性到高形式对称性。我们让对称操作作用在“维数 $\ge 1$ 的延展对象”上，例如威尔逊圈、't Hooft 圈以及其他更高维的推广。进一步地，可以把不同形式的对称性组合成一种称为“高群”（higher group）的代数结构。
  \item 超越群结构。我们不再用群、而用一种更弱的结构来描述对称性。可以逐步放松可逆性、结合律等要求，直到只剩下一个只能纳入高阶范畴论一般框架的非常复杂的结构。这类对称性称为非可逆对称性，或更一般地称为高阶范畴对称性。它们在物理与数学中都是活跃的研究方向\cite{Kong_2020,ji2020categorical,Johnson-Freyd:2020usu,Roumpedakis:2022aik,Bhardwaj:2022yxj,Bhardwaj:2022lsg,Bartsch:2022mpm,Bhardwaj:2022kot,Kaidi:2022cpf,Decoppet:2022dnz,Bhardwaj:2022maz,Bartsch:2022ytj,Delcamp:2023kew,Kaidi:2023maf,Bhardwaj:2023wzd,Bartsch:2023pzl,Bhardwaj:2023ayw,Bartsch:2023wvv,Decoppet:2023bay}。
\end{enumerate}

本讲义从不同视角介绍广义整体对称性。全文结构如下：在第~\ref{sec:ordinary}~节，我们在普通 $0$-形式对称性的语境下引入整体对称性的拓扑生成元框架。在第~\ref{sec:higher-form}~节，我们给出可逆高形式对称性的定义，并讨论纯麦克斯韦理论、带荷物质的麦克斯韦理论与非阿贝尔规范理论等例子。在第~\ref{sec:gauging}~节，我们讨论高形式对称性的规范化、't Hooft 反常，以及麦克斯韦理论与 3 维陈–西蒙斯理论的例子。在第~\ref{sec:three}~节，我们用三种等价的语言——拓扑缺陷网络、平坦联络与分类空间——对对称性生成元给出更全面的理解。这些语言对后文关于有限群对称性的讨论尤为重要。

在第~\ref{sec:anomaly-SPT}~节，我们用群上同调的语言介绍凝聚态物理中的对称保护拓扑（SPT）相与 't Hooft 反常，并给出 1+1 维 Haldane 链耦合 $0+1$ 维自旋-$\frac{1}{2}$ 边界以及 2+1 维 SPT 的例子。在第~\ref{sec:gen-SPT}~节，我们讨论 SPT 相的各种推广，包括高形式对称性与费米子系统，并给出用于分类的完整配边分类框架。

在第~\ref{sec:applications}~节，我们进一步讨论高形式对称性的应用。在第~\ref{sec:string}~节，我们讨论弦/M 理论框架下量子场论的几何工程，以及如何从额外维空间的拓扑计算高形式对称性；特别地，我们给出一个由 11 维 M 理论放在 Calabi–Yau 三流形（threefold）奇异点上得到 5 维 $\mc{N}=1$ 超共形场论的例子。在第~\ref{sec:appl-CMP}~节，我们讨论一些凝聚态物理应用，包括 2+1 维环面码与 3+1 维的高形式 SPT。

% ==== 原文 PDF 第 6 页 ====
最后，在第~\ref{sec:higher-group}~节，我们把视野扩展到更一般的范畴对称性，并特别强调高群对称性。我们在第~\ref{App:Cat}~节介绍范畴论的基础知识；在第~\ref{sec:2-group}~节用 2-范畴论的语言定义严格 2-群，并在第~\ref{sec:cross-extension}~节与第~\ref{sec:Postnikov}~节讨论它与弱 2-群的关系。我们在第~\ref{sec:2-group-phys}~节给出弱 2-群对称性的物理诠释，并在第~\ref{sec:2-group-gauging}~节讨论其规范化。弱 $n$-群与非可逆对称性分别在第~\ref{sec:n-group}~节和第~\ref{sec:non-invertible}~节作简要介绍。

在每一节开头，我们都会给出导读，帮助读者判断如何取舍并参与该节的内容。

\begin{yizhu}
此处“高群”指不同形式对称性混合成的代数结构：$0$-形式对称性与 $p$-形式对称性不再是直积，而是通过一个“混合反常”/Postnikov 类相互缠绕。这正对应后文第~5 节的 2-群：$H$（$1$-形式）被 $\Pi_1$（$0$-形式）以交叉模的方式作用，而 Postnikov 类 $\beta\in H^3_{\rm grp}(\Pi_1,\Pi_2)$ 度量这种缠绕的障碍。
\end{yizhu}

\begin{yizhu}
“非可逆对称性”一词的要点是：对称操作不再要求有逆，因而不再构成群，而构成融合范畴（甚至更高范畴）。此时对称性的“守恒律”以融合规则的形式出现，而不是以群元素的乘法表出现。本讲义只作一瞥（第~5.8 节），系统讨论可参见作者引用的综述。
\end{yizhu}

\begin{yizhu}[原文疑误]
本节末段（源码第 26 行）把 “symmetries” 误拼为 “symmtries”；此处按正确拼写理解并翻译，不影响内容。
\end{yizhu}
